\section{Тема 6. Линейный поиск}

\subsection{Постановка задачи (вариант~4)}

Задано $n$ линейных функций $y_1=k_1x+b_1,\,y_2=k_2x+b_2,\,\ldots,\,
y_n=k_nx+b_n$. Найти значение «верхней огибающей» этих функций
в заданной точке~$x$, то есть кусочно-линейной функции
\[
  f(x) = \max_{1 \le i \le n}\bigl(k_i x + b_i\bigr).
\]

\subsection{Математическое обоснование}

Поскольку все функции линейны, для фиксированного значения~$x$
вычисление верхней огибающей сводится к линейному поиску
максимума среди $n$ значений $\{k_i x + b_i\}$.
Сложность алгоритма~--- $O(n)$.

\subsection{Блок-схема алгоритма}

\begin{center}
\begin{tikzpicture}[
  node distance=10mm, start chain=going below,
  nodes={draw=black, top color=white, bottom color=lightgray!50,
         minimum width=5.4cm, align=center, font=\small},
  arrow/.style={->, >=Stealth, thick},
]
  \node[draw, rounded corners, on chain] {Start};
  \node[shape=trapezium, trapezium left angle=70,
        trapezium right angle=110, on chain, draw, join=by arrow]
       {$n$,\ $k_i,b_i$,\ $x$};
  \node[shape=rectangle, on chain, draw, join=by arrow]
       {$\mathrm{max}:=k_0 x{+}b_0,\ \mathrm{idx}:=0$};
  \node[shape=chamfered rectangle, chamfered rectangle xsep=10pt,
        on chain, draw, join=by arrow]
       (loop) {$i := 1,\ n{-}1$};
  \begin{scope}[local bounding box=scl]
    \node[shape=rectangle, on chain, draw, join=by arrow]
         {$v := k_i x{+}b_i$};
    \node[shape=diamond, aspect=2, on chain, draw, join=by arrow,
          minimum width=5.0cm]
         (c) {$v > \mathrm{max}$\,?};
    \node[shape=rectangle, draw, right=2cm of c]
         (upd) {$\mathrm{max}:=v;\ \mathrm{idx}:=i$};
    \draw[arrow] (c.east) -- node[above,font=\scriptsize]{$+$} (upd.west);
    \coordinate[on chain, join=by arrow] (le);
    \draw[arrow] (upd.south) |- (le);
  \end{scope}
  \draw[arrow] (scl.south) -| ($(scl.west)-(0.9,0)$) |- (loop);
  \coordinate[on chain] (sp);
  \draw[arrow] (loop) -| ($(scl.east)+(0.9,0)$) |- (sp);
  \node[shape=trapezium, trapezium left angle=70,
        trapezium right angle=110, on chain, draw, join=by arrow]
       {$\mathrm{max},\ \mathrm{idx}$};
  \node[draw, rounded corners, on chain, join=by arrow] {End};
\end{tikzpicture}
\end{center}

\subsection{Текст программы}

\lstinputlisting[style=Cstyle]{../src/t06.c}

\subsection{Тестовые данные}

\begin{center}
\begin{tabular}{l|l}
\toprule
Функции и точка & Огибающая \\
\midrule
$y=2x{+}1,\,y={-}x{+}4,\,y=0.5x{+}2$ при $x{=}1$ & $f(1)=3$ \\
\bottomrule
\end{tabular}
\end{center}

\textbf{Результат выполнения программы:}
\begin{lstlisting}[language={},basicstyle=\ttfamily\small,frame=single]
Введите количество функций n: 3
Введите коэффициенты k_i b_i для каждой функции:
2 1
-1 4
0.5 2
Введите значение x: 1
Верхняя огибающая в точке x=1.00: f(x)=3.0000 (функция 1)
\end{lstlisting}
